\scanpage{9}
figurer dans Bourbaki (pavages coniques, affines etc) -- tout au moins dans le cas \deleted{\unsure{alg}} non singulier. \inserted{Dans le cas infini,} les cas sembleraient liés (identiques ?) à celui où la forme quadratique invariante est (au signe près) de signature $(1,\unsure{--1})$ -- ce qui signifie que la sphère unité se décompose en deux composantes connexes \uline{simplement connexes} : c'est une de ces composantes qui est \inserted{doit} être le siège de la \deleted{\unsure{de la}} carte, on prend les traces dessus des cônes simpliciaux déduits (par op. de $G$) de l'un d'eux...

11) Lien entre l'icosaèdre et $E_8$ : la \deleted{\unsure{f}} forme quadratique de degré $4$ sur
\[
\mathbb Z[\alpha]=\mathbb Z[T]/(T^2+T-1)=\text{anneau des}
\]
entiers de $\mathbb Q(\tfrac12(-1\pm\sqrt5))=\mathbb Q(\sqrt5)$, composée avec la trace, donne la forme quadratique de $E_8$ .... d'où syst. de racines de $E_8$ .... Mais Attention : la ramification
\margin{?}

\mspage{8}
